Every figure below uses a placeholder file from \texttt{figures/}.
The class sets \verb|\graphicspath| to that folder, so \verb|\includegraphics| needs only the file name.

\subsection{Single-column figure}
This is the default case.
Size the image with \verb|width=\columnwidth| rather than a fixed length, so it adapts when the column model changes.
The label goes after the caption, never before, or the reference number comes out wrong.
See \cref{fig:single}.

\begin{figure}[t]
  \centering
  \includegraphics[width=\columnwidth]{placeholder-a}
  \caption{A single-column figure, sized with \texttt{width=\textbackslash columnwidth}.}
  \label{fig:single}
\end{figure}

\subsection{Full-width figure}
Use \verb|figure*| to span both columns in the two-column model.
It can only float to the top of a page, never to the middle or bottom, and never onto the page it was declared on.
Declare it a page or two before you want it to appear.
In the one-column model \verb|figure*| behaves exactly like \verb|figure|, so the source stays valid in both modes.
See \cref{fig:wide}.

\begin{figure*}[t]
  \centering
  
\includegraphics[width=0.8\textwidth]{placeholder-wide}
  \caption{A full-width figure, declared with \texttt{figure*} and sized against \texttt{\textbackslash textwidth}.}
  \label{fig:wide}
\end{figure*}

\subsection{Side-by-side subfigures}
\verb|subcaption| gives each panel its own letter and label.
Reference the whole figure as \cref{fig:pair}, or one panel as \cref{fig:pair-left}.
Keep the two \verb|subfigure| widths summing to slightly under 1, so the panels do not collide.

\begin{figure}[t]
  \centering
  \begin{subfigure}[b]{0.48\columnwidth}
    \includegraphics[width=\linewidth]{placeholder-a}
    \caption{First panel.}
    \label{fig:pair-left}
  \end{subfigure}
  \hfill
  \begin{subfigure}[b]{0.48\columnwidth}
    \includegraphics[width=\linewidth]{placeholder-b}
    \caption{Second panel.}
    \label{fig:pair-right}
  \end{subfigure}
  \caption{Two subfigures side by side.}
  \label{fig:pair}
\end{figure}

\subsection{Stacked subfigures}
Drop the \verb|\hfill| and add a blank-line break to stack panels vertically instead.

\begin{figure}[t]
  \centering
  \begin{subfigure}[b]{0.7\columnwidth}
    \includegraphics[width=\linewidth]{placeholder-a}
    \caption{Top panel.}
    \label{fig:stack-top}
  \end{subfigure}

  \vspace{0.6em}

  \begin{subfigure}[b]{0.7\columnwidth}
    \includegraphics[width=\linewidth]{placeholder-b}
    \caption{Bottom panel.}
    \label{fig:stack-bottom}
  \end{subfigure}
  \caption{Two subfigures stacked vertically.}
  \label{fig:stack}
\end{figure}

\subsection{Text-wrapped figure}
\verb|wrapfig| lets prose flow around a small figure.
It is fragile near page breaks, list environments, and section headings, so use it sparingly and check the output every time.
Put it at the start of a paragraph, never inside one.

\begin{wrapfigure}{r}{0.45\columnwidth}
  \centering
  \includegraphics[width=0.43\columnwidth]{placeholder-b}
  \caption{A wrapped figure.}
  \label{fig:wrapped}
\end{wrapfigure}

The paragraph that follows the \verb|wrapfigure| environment is the one that wraps around it.
Give the environment a width a little larger than the image itself, so the text does not press against the edge of the graphic.
If the wrap breaks across a page, move the environment one paragraph earlier or later rather than fighting the placement options.
Two or three sentences of surrounding prose are usually enough to fill the space beside a half-column image.

\subsection{Rotated and oversized figures}
\verb|adjustbox| handles an image that refuses to fit.
\verb|max width=\columnwidth| shrinks only when needed, and leaves smaller images untouched.

\begin{figure}[t]
  \centering
  \begin{adjustbox}{max width=\columnwidth}
    \includegraphics[width=1.6\columnwidth]{placeholder-a}
  \end{adjustbox}
  \caption{An oversized image, clamped by \texttt{adjustbox}.}
  \label{fig:clamped}
\end{figure}
